% Full connective/rule surface of Pfenning's two-context ILL sequent system,
% Handout 4, draft 27 January 1998, Section 2.6, printed pages 42--45.
% This is not the inherited four-rule bang extension or one-sided CLL.
\calculussheet{sheet:p14-ill-pfenning}{Intuitionistic Linear Logic: Pfenning's Two-Context Presentation}{Language: first-order predicate atoms over a fixed term signature, multiplicative $1,\otimes,\multimap$, additive $\top,0,\mathbin{\&},\oplus$, exponential $!$, unrestricted implication $\supset$, and $\forall,\exists$; judgments: $\Gamma;\Delta\Rightarrow A$ with one conclusion; $\Gamma$ contains reusable labelled assumptions, $\Delta$ linear labelled occurrences; exchange is implicit, not linear weakening or contraction; negation and equality: neither primitive.}

\noindent\emph{Source:} \cite[\S2.6, pp.~42--45]{Pfenning1998ILLSequent}.
Labels on assumptions are suppressed in the displayed source abbreviation,
not identified.  The dot denotes an empty context; $\Delta_1\times\Delta_2$
denotes combination of disjoint labelled linear contexts, not a Cartesian
product of sets.  The unrestricted context is retained in both premises.
All substitutions below avoid capture; eigenparameters are fresh for the
conclusion and its undischarged hypotheses.

\hbrule{rule:p14-ill-pf-id}{I}{}{\Gamma;A\Rightarrow A}{the linear context consists of exactly the displayed occurrence}
\hbrule{rule:p14-ill-pf-copy}{DL}{(\Gamma,A);(\Delta,A)\Rightarrow C}{(\Gamma,A);\Delta\Rightarrow C}{the premise uses a fresh linear copy of the displayed unrestricted assumption; the unrestricted occurrence remains available}
\hbrule{rule:p14-ill-pf-lolli-r}{$\multimap R$}{\Gamma;\Delta,A\Rightarrow B}{\Gamma;\Delta\Rightarrow A\multimap B}{\none}
\hbrule{rule:p14-ill-pf-lolli-l}{$\multimap L$}{\Gamma;\Delta_1\Rightarrow A\qquad\Gamma;\Delta_2,B\Rightarrow C}{\Gamma;\Delta_1\times\Delta_2,A\multimap B\Rightarrow C}{the linear premise contexts are disjoint}
\hbrule{rule:p14-ill-pf-tensor-r}{$\otimes R$}{\Gamma;\Delta_1\Rightarrow A\qquad\Gamma;\Delta_2\Rightarrow B}{\Gamma;\Delta_1\times\Delta_2\Rightarrow A\otimes B}{the linear premise contexts are disjoint}
\hbrule{rule:p14-ill-pf-tensor-l}{$\otimes L$}{\Gamma;\Delta,A,B\Rightarrow C}{\Gamma;\Delta,A\otimes B\Rightarrow C}{the two displayed premise occurrences replace the principal tensor occurrence}
\hbrule{rule:p14-ill-pf-one-r}{$1R$}{}{\Gamma;\cdot\Rightarrow1}{the linear context is empty}
\hbrule{rule:p14-ill-pf-one-l}{$1L$}{\Gamma;\Delta\Rightarrow C}{\Gamma;\Delta,1\Rightarrow C}{\none}
\hbrule{rule:p14-ill-pf-with-r}{$\&R$}{\Gamma;\Delta\Rightarrow A\qquad\Gamma;\Delta\Rightarrow B}{\Gamma;\Delta\Rightarrow A\mathbin{\&}B}{the same linear context occurs in the two alternative proofs}
\hbrule{rule:p14-ill-pf-with-l1}{$\&L_1$}{\Gamma;\Delta,A\Rightarrow C}{\Gamma;\Delta,A\mathbin{\&}B\Rightarrow C}{\none}
\hbrule{rule:p14-ill-pf-with-l2}{$\&L_2$}{\Gamma;\Delta,B\Rightarrow C}{\Gamma;\Delta,A\mathbin{\&}B\Rightarrow C}{\none}
\hbrule{rule:p14-ill-pf-top-r}{$\top R$}{}{\Gamma;\Delta\Rightarrow\top}{the displayed linear context may be arbitrary}
\hbrule{rule:p14-ill-pf-plus-r1}{$\oplus R_1$}{\Gamma;\Delta\Rightarrow A}{\Gamma;\Delta\Rightarrow A\oplus B}{\none}
\hbrule{rule:p14-ill-pf-plus-r2}{$\oplus R_2$}{\Gamma;\Delta\Rightarrow B}{\Gamma;\Delta\Rightarrow A\oplus B}{\none}
\hbrule{rule:p14-ill-pf-plus-l}{$\oplus L$}{\Gamma;\Delta,A\Rightarrow C\qquad\Gamma;\Delta,B\Rightarrow C}{\Gamma;\Delta,A\oplus B\Rightarrow C}{the same residual linear context occurs in both alternatives}
\hbrule{rule:p14-ill-pf-zero-l}{$0L$}{}{\Gamma;\Delta,0\Rightarrow C}{\none}
\hbrule{rule:p14-ill-pf-forall-r}{$\forall R^a$}{\Gamma;\Delta\Rightarrow[a/x]A}{\Gamma;\Delta\Rightarrow\forall x.A}{$a$ is a fresh eigenparameter, absent from the conclusion sequent}
\hbrule{rule:p14-ill-pf-forall-l}{$\forall L$}{\Gamma;\Delta,[t/x]A\Rightarrow C}{\Gamma;\Delta,\forall x.A\Rightarrow C}{$t$ is a term in the current parameter signature}
\hbrule{rule:p14-ill-pf-exists-r}{$\exists R$}{\Gamma;\Delta\Rightarrow[t/x]A}{\Gamma;\Delta\Rightarrow\exists x.A}{$t$ is a term in the current parameter signature}
\hbrule{rule:p14-ill-pf-exists-l}{$\exists L^a$}{\Gamma;\Delta,[a/x]A\Rightarrow C}{\Gamma;\Delta,\exists x.A\Rightarrow C}{$a$ is a fresh eigenparameter, absent from the conclusion sequent}
\hbrule{rule:p14-ill-pf-implies-r}{$\supset R$}{(\Gamma,A);\Delta\Rightarrow B}{\Gamma;\Delta\Rightarrow A\supset B}{the new antecedent occurrence is unrestricted}
\hbrule{rule:p14-ill-pf-implies-l}{$\supset L$}{\Gamma;\cdot\Rightarrow A\qquad\Gamma;\Delta,B\Rightarrow C}{\Gamma;\Delta,A\supset B\Rightarrow C}{the first premise has no linear assumptions}
\hbrule{rule:p14-ill-pf-bang-r}{$!R$}{\Gamma;\cdot\Rightarrow A}{\Gamma;\cdot\Rightarrow!A}{the linear context is empty in premise and conclusion}
\hbrule{rule:p14-ill-pf-bang-l}{$!L$}{(\Gamma,A);\Delta\Rightarrow C}{\Gamma;\Delta,!A\Rightarrow C}{the premise moves the contents of the principal bang to the unrestricted context}

There is no $\top$ left rule, no $0$ right rule, and no primitive Cut in
this cut-free presentation.  The source's Theorems~2.9--2.10 compare this
system with its own normal/atomic natural deductions; they are not a
transfer of the one-sided CLL phase or cut theorem to this sheet.
